\section{Estimation: AIPW, cross-fitting, and which covariates to adjust for}
\label{sec:estimation}

\emph{Source: handout \S1.2--1.3. Status: skeleton from a first read.}

\subsection{Estimating one level, discrete \texorpdfstring{$A$}{A} (\S1.2, mostly REUSE)}

Three classical estimators (Robins/Rotnitzky/Zhao 1994; Hahn 1998):
regression (regress $Y^*$ on $(A,Z)$, average), inverse-propensity
weighting, or \textbf{AIPW}, which combines both and is \emph{doubly
robust} --- consistent if \emph{either} piece is right:
\begin{equation}
  \hat\theta(a) \;=\; \frac{1}{n}\sum_{i=1}^n \left[
    \frac{\mathbf{1}(A_i=a)}{\hat\pi(a\mid Z_i)}\big(Y_i^* - \hat\mu(a,Z_i)\big)
    + \hat\mu(a,Z_i) \right].
  \label{eq:aipw}
\end{equation}
\emph{In words, person by person:} start with the regression's best guess
for person $i$'s outcome if they'd had exercise level $a$, given their
covariates ($\hat\mu(a,Z_i)$ --- this term alone is used for everyone). Then,
\emph{only for the people who actually did exercise at level $a$}
($\mathbf{1}(A_i=a)$ switches this next part on or off), add a correction:
how far their real outcome was from that regression guess
($Y_i^* - \hat\mu(a,Z_i)$), inflated by how rare it was for someone with
their covariates to end up at exercise level $a$
($1/\hat\pi(a\mid Z_i)$ --- rarer means a bigger correction, since that
person is doing more work standing in for others like them). Average this
person-by-person quantity over everyone, and that's $\hat\theta(a)$. The
``doubly robust'' part: if $\hat\mu$ is wrong but $\hat\pi$ is right, the
correction term fixes it on average; if $\hat\pi$ is wrong but $\hat\mu$ is
right, the correction term is small anyway because the regression guess was
already good.

\textbf{Cross-fitting} \reusetag{Chernozhukov et al.\ 2018}: fit
$\hat\mu, \hat\pi$ on one data half, evaluate on the other, swap. This lets
flexible ML be used for the nuisance functions while still getting a valid
Wald interval $\hat\theta(a) \pm 1.96\,\widehat{\mathrm{se}}$, \emph{provided}
the two nuisance errors shrink fast enough \emph{jointly} (their product
must be $o_p(n^{-1/2})$). Cross-fitting removes a technical condition, not
the underlying accuracy requirement --- a common misreading to watch for.

\subsection{Which covariates \texorpdfstring{$Z$}{Z} to adjust for (\S1.3)}

\reusetag{Rotnitzky \& Smucler 2020}, \newtag~application. The parents $W$
are a \emph{valid} adjustment set but not necessarily the lowest-variance
one. Adding non-descendants of $A$ that help predict $Y^*$ lowers variance
without changing the answer (same logic as good covariates in a randomized
trial).

\textbf{Key result:} when $f$ uses all features, the set of \emph{all
non-descendants of $A$} is never worse than any other valid set in large
samples (among the adjustment estimators used here). Two practical facts:
\begin{enumerate}[label=(\alph*)]
  \item The \emph{population} positivity condition does not get harder ---
    because $P(a \mid \text{non-descendants}) = P(a \mid W)$ by the graph.
    What gets harder is \emph{fitting} a bigger regression in finite
    samples. So: \textbf{fit the propensity on $W$, fit the outcome
    regression on the larger set.}
  \item This result uses only the graph's \emph{arrows} --- never its
    equations. (The handout notes this was verified exhaustively on small
    graphs in the review --- a good sanity-check exercise to redo.)
\end{enumerate}

\begin{openquestions}
  \item Work through \emph{why} adding non-descendants can only help (not
    hurt) asymptotic variance --- is there a clean intuition beyond ``more
    relevant predictors of $Y^*$ reduce residual variance''?
  \item Handout \S1.9 marks ``cross-fitted AIPW at discrete levels; coverage
    on oracle testbed'' as intern-feasible / moderate supervision
    (\feasmod), and ``parents vs.\ efficient set (variance drop)'' likewise
    --- plausible first coding tasks once toy-SCM verification
    (\S\ref{sec:identification}) is done.
\end{openquestions}
